The Massey forensic data, the two biologically untenable conditions, and the critical threshold from bias to fabrication

The forensic re-analysis of the public raw sequencing reads associated with RaTG13 constitutes a decisive empirical correction to one of the longest-standing anchors of pure natural-proximal narratives. For more than a decade the sample has been presented, in publications and public communications supporting the zoonosis hypothesis, as a clean environmental fecal swab collected from a Rhinolophus bat in the Mojiang/Tongguan mine in 2013. That characterization supplied the strongest claimed natural, wild-origin data point linking a close relative (~96.2 % identity) to a typical bat-reservoir sampling method. The Massey forensic examination of those same public raw reads demonstrates that the metagenomic profile is inconsistent with a typical fecal swab: bacterial content is anomalously low, while eukaryotic and transcriptome-derived signals are disproportionately high. The profile aligns instead with material of reproductive-tract or mating-plug origin. The discrepancy is not a minor technical quibble; it removes the prior strongest claim of a clean wild 2013 environmental fecal-swab anchor and thereby further erodes the empirical foundation for an undetected natural proximal ancestor already carrying the exact PRRA + CGG-CGG furin-cleavage configuration.

This correction does not stand alone. It intersects with the broader documented absence, detailed in the public record of the Year-5 RPPR for NIH Grant R01AI110964, of any natural intermediate-host specimen bearing the precise polybasic site. Intensive sampling through 31 May 2019 recovered diverse SARS-related sequences yet never recovered the defining genomic feature of the virus that first became detectable in Wuhan later that year. The Massey result simply eliminates one of the few remaining candidate anchors that had been repeatedly cited as evidence of ordinary wildlife sampling.

Against this backdrop the practical possibility of recovering a temporally authentic 2019 natural proximal ancestor from any living reservoir in 2026 has reached zero. That zero is reinforced by two independent biological constraints, neither of which is tenable.

The first constraint is the continuous operation of the endogenous molecular clock. RNA viruses accumulate substitutions at measurable rates while they replicate. Any lineage circulating in 2019 and still replicating through 2026 has experienced six to seven additional years of mutational accumulation. A contemporary specimen is therefore a descendant generation; it cannot be the original 2019 genome. For the lineage to remain recoverable as the precise proximal ancestor it would have to have frozen its clock—ceased ordinary evolution—for more than half a decade. No mechanism in RNA-virus biology permits replication to continue while substitution, recombination and drift are suspended. A lineage cannot freeze its clock for half a decade.

The second constraint is pharmacological and population-level. After the clinical deployment of molnupiravir, NHC-TP incorporation elevated mutation rates and imposed a directional bias toward G-to-A and C-to-U transitions. These signature genomes have been recovered in global databases, appear as long phylogenetic branches, and in multiple documented instances have seeded viable descendant clades that continue to circulate. The mutational load has therefore crossed from treated individuals into the broader population. For a hypothetical 2019 proximal ancestor to remain intact in 2026 its descendants would have had to avoid every transmission chain exposed to this pressure or, having acquired the excess transitions, would have had to reverse them with perfect precision while restoring the original codon-usage bias and phenotypic state. Global avoidance is statistically implausible; perfect reversal is chemically and population-genetically impossible. A lineage cannot globally avoid or perfectly reverse the cross-population NHC-TP mutational load.

Neither condition is biologically tenable. Their joint action permanently closes the recovery window. Future sampling may yield closer overall relatives or independent furin-like motifs in distant lineages; it cannot yield a genome that both stopped evolving in 2019 and remained insulated from, or perfectly erased, the directional mutational process that has been circulating since 2022.

The persistence of narratives that continue to treat a pristine 2019 natural proximal ancestor as recoverable, or that continue to present RaTG13 as an unproblematic wild fecal swab despite the Massey forensic profile, can be examined through the progressive sequence of asymmetric judgment. Asymmetric judgment does not begin as deliberate falsehood. It originates as a cognitive bias that privileges one party—the preferred natural-proximal account—while systematically disadvantaging the alternative. Under sustained pressure to maintain that preferred narrative the bias advances through three recognizable stages until it crosses a definitive threshold into conscious deception.

In the first stage, biased interpretation defends the favorite through selective emphasis. The fecal-swab characterization of RaTG13 is retained because it supplies emotional and narrative comfort; the Massey metagenomic discrepancies are minimized or contextualized as technical noise. Contradictory sampling data from the intensive pre-2019 surveys are filtered so that the absence of the exact FCS-bearing specimen is framed as temporary incompleteness rather than structural negative evidence. The evaluator still operates within perception rather than invention; the complete record is not yet severed.

In the second stage, evidence is altered or selectively omitted to preserve the preferred illusion. The forensic profile inconsistent with a fecal swab is left out of subsequent summaries; the irreversible advance of both the endogenous clock and the NHC-TP mutational load is omitted from discussions of future sampling prospects. Timelines are adjusted so that the six-to-seven-year gap appears bridgeable; the cross-population seeding of signature clades is quietly set aside. The resulting account retains a superficial continuity with known events yet has been surgically edited to eliminate the two untenable conditions.

In the third stage, reality itself is invented to rescue a crumbling narrative. When selective omission proves insufficient, external explanations are constructed: the Massey signals are attributed to unspecified contamination or laboratory artifact without supporting evidence; the pharmacological mutational load is declared irrelevant to natural reservoirs; or entirely new, still-unsampled colonies are posited that somehow froze their clocks and avoided every NHC-TP exposure. At this juncture the defense has ceased to be perspective or even selective presentation; it has become an active substitution of preferred story for observed reality.

The transition crystallizes into absolute falsehood at the precise moment the evaluator knows that the original benchmarks—the fecal-swab claim and the recoverability of a 2019 proximal ancestor—have been violated by the Massey data and by the two biological constraints, yet continues to assert that those standards remain intact. Beyond this threshold, asymmetric judgment is no longer a cognitive distortion amenable to correction by additional information. It is a deliberate construction of untruth engineered to maintain the fragile hierarchy of the preferred narrative. Recognition of the sequence supplies the analytical clarity required to identify when the continued defense of a recoverable pristine 2019 natural specimen, or of an unproblematic fecal-swab anchor, has crossed from contested interpretation into constructed deception.
